\subsection{Out-of-domain stress test}
\label{sec:ood}

RefCOCO-family benchmarks are photo datasets with COCO-style objects.
They do not establish generalization to interfaces, documents, robotics images, microscopy, or aerial imagery.
We define a separate qualitative stress test with 10 to 30 images per domain, one hand-written referring expression per image, and human scoring.
No OOD numbers are claimed in this draft.

\begin{table}[H]
  \centering
  \caption{OOD protocol. Prompts are examples; the image set has not been scored.}
  \label{tab:ood-domains}
  \scriptsize
  \begin{tabularx}{\textwidth}{@{}l c X l@{}}
    \toprule
    Domain & Target $N$ & Example prompt & Status \\
    \midrule
    UI         & 10 to 30 & blue submit button & images missing \\
    Robotics   & 10 to 30 & bottle near gripper & images missing \\
    Aerial     & 10 to 30 & building by the road & images missing \\
    Microscopy & 10 to 30 & dark region at edge & images missing \\
    Documents  & 10 to 30 & title block at top & images missing \\
    \bottomrule
  \end{tabularx}
\end{table}

For each image, the planned comparison runs Locate-SAM2 hybrid and DINO-Tiny + SAM2 with the same adapter configuration as Table~\ref{tab:main}.
Reviewers score four binary criteria: prompt understood, correct instance, usable box, and usable mask.
The reported metric will be the percentage of examples with a usable final mask.
All RefCOCO-family validation results reported in this paper are complete.
OOD results remain outside the claims until image sets and human scores are present under \texttt{experiments/ood/}.
